\begin{exsol}{12}{Cabral 3.14 Intel}{63}{Exam}{A2}
\asked Explain Intel's sustained competitive advantage over its rivals and weigh the
explanatory power of the four causes of \S3.5: impediments to imitation, causal ambiguity,
strategy and history.

\method Sources of sustained competitive advantage (Cabral \S3.5, pp.~56--59): impediments to
imitation (legal and other, p.~57), causal ambiguity and tacit knowledge (p.~57), strategy
(p.~57), history and the learning curve (pp.~58--59). Facts about Intel are restricted to what
Cabral reports: leadership in quality and efficient production, very high fab set-up costs and
a much smaller rival AMD (\S5.3, pp.~103--104); exclusivity rebates and the resulting antitrust
cases (\S12.2, pp.~315--316; \S12.4, pp.~321--322); reliance on production efficiency and the
learning curve rather than on intellectual property (\S15.3, p.~386).

\solution
\subparagraph{(a) Impediments to imitation.} Two barriers apply. The first is scale: the set-up
cost of a fab is so high that a rival can hardly match Intel's size (p.~103), so unit costs stay
lower for the leader even when the technology is known. This is less ``imitation is
impossible'' than ``imitation at the same scale is unprofitable''. The second is legal
protection: patents and secrecy exist, but Cabral notes that Intel relies primarily on production
efficiency and that its success ``is more a case of execution than building valuable IP assets''
(p.~386). Unlike Xerox's plain-paper patent (p.~57),
law is not the main barrier. \emph{Weight: high for scale, low for patents.}

\subparagraph{(b) Causal ambiguity.} Producing chips efficiently means running a fab with high
yields at each new generation, a body of process know-how built by experience and held by teams
as tacit knowledge (p.~57). A rival that hires some of Intel's engineers does not acquire the
system, much as Toyota's procurement system could not be copied piece by piece (p.~57).
\emph{Weight: moderate to high}, because it explains why the efficiency lead persists rather
than spreading.

\subparagraph{(c) Strategy.} Cabral documents an aggressive strategy: all-unit discounts that
made buying any AMD chips expensive for PC makers, an instance of raising rivals' costs (\S12.2,
pp.~315--316), which led to large antitrust settlements and fines (\S12.4, pp.~321--322); he
suggests aggressive pricing is one reason AMD did not grow into a closer competitor (p.~103).
Moving quickly from one generation to the next (p.~386) is also a strategic choice: it keeps the
leader ahead on the learning curve of each new product. \emph{Weight: high}, though part of it
(the rebates) was later restricted by antitrust.

\subparagraph{(d) History.} The learning-curve mechanism converts an early lead into a lasting
one: the firm that has produced more has lower costs, which lets it price lower and produce more
again, as in Boeing's case (pp.~58--59). An early dominant position in microprocessors (Cabral
lists Intel as the dominant firm since the 1980s, p.~104) also matters because buyers and
complementary products (software, operating systems) are built around the dominant design
(complementary products are the setting of Exercise 13.7). \emph{Weight: high}, as the starting
condition the other three amplify.

\subparagraph{(e) Weighing them.} The causes are complements, not rivals. History gave an early
lead; the learning curve and fab scale (history plus a scale-based impediment) made the lead
cumulative; tacit process know-how (causal ambiguity) kept rivals from closing the gap by
copying; and strategy (pricing, rebates, fast generations) protected the share the cost
advantage earned. In Cabral's own reading the core is production efficiency along the learning
curve (p.~386), so the ranking is: scale and learning (history) and strategy first, causal
ambiguity second, legal protection last.

\interp For the course, the useful lesson is the one Aula~1 draws from Toyota: an advantage that
rests on a coherent system of mutually reinforcing elements (scale, know-how, pricing, product
cadence) is harder to imitate than one resting on a single asset, because a rival must copy the
whole system at once (\MR{} Ch.~4 complementarities; Cabral p.~57). Sustained is not permanent:
the same framework says the advantage erodes if the learning-curve lead is lost, as with
Lockheed and McDonnell Douglas against Boeing (p.~59).

\verify Consistency check with \S3.5's evidence: performance differences within an industry are
large and persistent, and only about 20\% of the variance in profit rates is explained by
industry and size (p.~56). A single dominant firm whose lead persists across generations is
what the four causes predict, and the antitrust record (pp.~321--322) shows that at least part
of the advantage was defended by conduct, not only earned by efficiency.

\answer{Intel's advantage is best explained by history and the learning curve combined with fab
scale economies (an early lead that became a cost lead), protected by strategy (aggressive
pricing, exclusivity rebates, fast product generations) and made hard to copy by tacit process
know-how (causal ambiguity). Legal protection (patents) is the weakest of the four in Cabral's
account (p.~386). The causes reinforce each other, which is why the advantage persisted.}
\end{exsol}
